\documentclass[12pt, a4paper]{article}

\usepackage[utf8]{inputenc}
\usepackage[T1]{fontenc}

\usepackage{textcomp}
\usepackage{geometry}
\geometry{margin=2.5cm}
\usepackage{titlesec}
\usepackage{enumitem}
\usepackage{amsmath, amssymb}
\usepackage{hyperref}
\usepackage{xcolor}
\usepackage{mdframed}


% Colori
\definecolor{accent}{RGB}{30, 80, 140}
\definecolor{boxbg}{RGB}{240, 245, 255}
\definecolor{fontecolor}{RGB}{90, 90, 90}


% Box fonti
\newmdenv[
backgroundcolor=boxbg,
linecolor=accent,
linewidth=1pt,
leftmargin=0pt,
rightmargin=0pt,
innerleftmargin=10pt,
innerrightmargin=10pt,
innertopmargin=6pt,
innerbottommargin=6pt,
skipabove=8pt,
skipbelow=8pt
]{fontebox}

\newcommand{\fonti}[1]{%
	\begin{fontebox}
		\textcolor{fontecolor}{\footnotesize\textbf{Fonti:} #1}
	\end{fontebox}%
}

% Comando per nota concettuale
\newcommand{\nota}[1]{%
	\par\smallskip
	\noindent\textit{\small $\triangleright$\; #1}
	\par\smallskip
}

\setlist[itemize,1]{leftmargin=1.5em, itemsep=2pt, topsep=4pt}
\setlist[itemize,2]{leftmargin=1.5em, itemsep=1pt, topsep=2pt}

\hypersetup{
	colorlinks=true,
	linkcolor=accent,
	urlcolor=accent
}

\title{%
	\textbf{Strutture Hamiltoniane in Meccanica Quantistica}\\[6pt]
	\normalsize Indice Concettuale
}
\author{Tommaso Pedroni}
\date{\today}

\begin{document}
	
	\maketitle
	\tableofcontents
	\newpage
	
	% ============================================================
	\section*{Domanda Guida della Tesi}
	\addcontentsline{toc}{section}{Domanda Guida della Tesi}
	% ============================================================
	
	Le strutture della meccanica razionale --- spazio delle fasi, parentesi di Poisson,
	flusso hamiltoniano, teorema di Noether --- sono letteralmente presenti nella meccanica
	quantistica. La tesi argomenta che questo non è una coincidenza né un'analogia formale,
	ma una conseguenza strutturale leggibile in due modi complementari.
	
	\medskip
	
	\textbf{Bottom-up (Parti I--II + Cap.\ 4).} Meccanica razionale e meccanica quantistica
	sono entrambe $C^*$-algebre --- una commutativa, l'altra no --- connesse da un campo
	continuo tramite la strict deformation quantization di Rieffel. Il Lie-bracket classico
	e il commutatore quantistico sono la stessa struttura astratta in due rappresentazioni
	distinte.
	
	\medskip
	
	\textbf{Top-down (Cap.\ 3).} Lo spazio degli stati quantistico $\mathbb{P}\mathcal{H}$
	è una varietà simplettica kähleriana. Forma simplettica, campi hamiltoniani, mappa
	momento e teorema di Noether non sono importati dall'esterno: emergono dalla
	decomposizione del prodotto scalare di $\mathcal{H}$.
	
	\medskip
	
	La risposta si articola in quattro passi con dipendenza logica stretta.
	
	\medskip
	
	\textbf{Passo 1 (Cap.\ 1).} La MR è costruita in forma algebrica invariante e
	riconosciuta come $C^*$-algebra commutativa tramite Gelfand. Il Lie-bracket classico
	è struttura interna di $C^\infty_c(M)\subset C_0(M)$, non un assioma aggiunto.
	
	\medskip
	
	\textbf{Passo 2 (Cap.\ 2).} La MQ è costruita senza postulare $\mathcal{H}$.
	L'ostruzione di Wintner (algebre di Banach) motiva l'algebra di Weyl. Lo spazio
	di Hilbert emerge dal teorema GNS. Il commutatore $[a,b]=ab-ba$ è struttura interna
	di ogni $C^*$-algebra non commutativa --- parallelo esatto del Passo~1.
	
	\medskip
	
	\textbf{Passo 3 (Cap.\ 3).} $\mathbb{P}\mathcal{H}$ è una varietà simplettica
	kähleriana. L'equazione di Schrödinger, il teorema di Noether e le relazioni di
	indeterminazione sono teoremi su questa varietà, non assiomi indipendenti.
	
	\medskip
	
	\textbf{Passo 4 (Cap.\ 4).} La strict deformation quantization di Rieffel costruisce
	una famiglia continua $\{A_\hbar\}_{\hbar\in[0,1]}$ di $C^*$-algebre. La condizione
	di Dirac garantisce che parentesi di Poisson e commutatore siano la stessa struttura
	in senso normato per $\hbar\to 0$.
	
	% ============================================================
	\newpage
	\part*{Parte I --- Il Lato Classico}
	\addcontentsline{toc}{part}{Parte I: Il Lato Classico}
	% ============================================================
	
	\nota{La Parte~I costruisce il linguaggio geometrico e algebrico della meccanica
		razionale in forma invariante, poi riconosce la struttura risultante come
		$C^*$-algebra commutativa. Ogni definizione introdotta qui ha un analogo esplicito
		nel Cap.\ 3.}
	
	% ============================================================
	\section{Varietà Simplettiche e Algebra delle Osservabili}
	% ============================================================
	
	\nota{Sequenza logica: algebre di Lie forniscono il linguaggio del bracket; la geometria
		simplettica introduce $C^\infty(M)$ come spazio naturale delle osservabili; le
		parentesi di Poisson dotano $C^\infty(M)$ di struttura di algebra di Lie; il teorema
		di Gelfand riconosce $C_0(M)$ come $C^*$-algebra. Il capitolo si chiude con un
		risultato: la MR \emph{è} la teoria delle $C^*$-algebre commutative.}
	
	% ----------------------------------------------------------
	\subsection{Algebre di Lie}
	% ----------------------------------------------------------
	
	\fonti{Humphreys, \textit{Introduction to Lie Algebras and Representation Theory};
		Kirillov, \textit{An Introduction to Lie Groups and Lie Algebras}.}
	
	\nota{Prerequisito per l'intero capitolo. Le algebre di Lie sono il linguaggio in cui
		si esprimono sia le parentesi di Poisson (\S1.3) sia il commutatore quantistico
		(Cap.\ 2). Ogni bracket che compare in questa tesi --- classico o quantistico ---
		è una rappresentazione della struttura astratta definita qui.}
	
	\begin{itemize}
		\item \textbf{Definizione assiomatica} $(\mathfrak{g},[\cdot,\cdot])$: bilinearità,
		antisimmetria, identità di Jacobi. L'identità di Jacobi è anticipata come
		condizione di coerenza strutturale --- il suo legame con $d\omega=0$ è il punto
		tecnico cruciale di \S1.3.
		\item \textbf{Omomorfismi, rappresentazioni} e algebra inviluppante universale
		(enunciato).
		\item \textbf{Esempi rilevanti:} $\mathfrak{su}(2)$, $\mathfrak{u}(n)$; mappa
		esponenziale $\exp:\mathfrak{g}\to G$.
	\end{itemize}
	
	% ----------------------------------------------------------
	\subsection{Geometria Simplettica}
	% ----------------------------------------------------------
	
	\fonti{Abraham \& Marsden, \textit{Foundations of Mechanics};
		Cannas da Silva, \textit{Lectures on Symplectic Geometry}.}
	
	\nota{Qui entra $C^\infty(M)$: è lo spazio in cui le hamiltoniane vivono naturalmente,
		perché definire $X_f$ tramite $\iota_{X_f}\omega = df$ richiede differenziabilità
		di $f$. Il passaggio a $C_0(M)$ avverrà in \S1.5, quando si vorrà una struttura
		di $C^*$-algebra.}
	
	\begin{itemize}
		\item \textbf{Nota sulla compattezza:} salvo diversa indicazione, $M$ è una
		varietà simplettica localmente compatta di Hausdorff. I risultati che richiedono
		compattezza di $M$ sono segnalati esplicitamente con \emph{(M compatta)}.
		\item \textbf{Varietà simplettica} $(M,\omega)$: $\omega$ chiusa ($d\omega=0$) e
		non degenere; dimensione pari, orientazione canonica.
		\item \textbf{Isomorfismi musicali} $\flat:TM\to T^*M$ e $\sharp:T^*M\to TM$
		indotti da $\omega$; campi vettoriali hamiltoniani $X_f$ via $\iota_{X_f}\omega=df$;
		esistenza e unicità dall'isomorfismo musicale.
		\item \textbf{Teorema di Darboux:} esistenza \emph{locale} di coordinate
		$(q^i,p_i)$ con $\omega=\sum_i dq^i\wedge dp_i$. Le coordinate canoniche sono
		una conseguenza locale --- non esistono coordinate di Darboux globali su varietà
		compatte come $S^2$.
		\item \textbf{Simplettomorfismi e teorema di Liouville:} il flusso di ogni campo
		hamiltoniano è un gruppo a un parametro di simplettomorfismi \emph{(completezza
			garantita se $M$ è compatta, o con ipotesi di crescita su $H$)}; invarianza
		della misura di Liouville $\omega^n/n!$.
	\end{itemize}
	
	% ----------------------------------------------------------
	\subsection{Parentesi di Poisson e Struttura di Algebra di Lie}
	% ----------------------------------------------------------
	
	\fonti{Abraham \& Marsden, \textit{Foundations of Mechanics};
		Cannas da Silva, \textit{Lectures on Symplectic Geometry}.}
	
	\begin{itemize}
		\item \textbf{Parentesi di Poisson:} $\{f,g\}=\omega(X_f,X_g)$; bilinearità,
		antisimmetria, regola di Leibniz.
		\item \textbf{Identità di Jacobi $\Leftrightarrow$ $d\omega=0$:} l'equivalenza
		si dimostra espandendo $\{\{f,g\},h\}$ tramite la formula di Cartan
		$\mathcal{L}_X=\iota_X d+d\iota_X$ e usando la chiusura di $\omega$. Riferimento:
		Abraham--Marsden, Thm~2.4.6; Cannas da Silva, Prop.~2.3. Questo è il punto
		tecnico cruciale: la stessa equivalenza garantirà Jacobi su $\mathbb{P}\mathcal{H}$
		nel Cap.~3.
		\item \textbf{$(C^\infty(M),\{\cdot,\cdot\})$ come algebra di Lie:} tutti gli
		assiomi di \S1.1 soddisfatti; il Lie-bracket classico è la struttura astratta
		di cui il commutatore $\frac{1}{i\hbar}[\cdot,\cdot]$ è l'omologo non commutativo.
		\item \textbf{Mappa $f\mapsto X_f$:} omomorfismo di algebre di Lie
		$(C^\infty(M),\{\cdot,\cdot\})\to(\mathfrak{X}(M),[\cdot,\cdot])$; nucleo nelle
		funzioni costanti.
		\item \textbf{Evoluzione delle osservabili:} $\tfrac{d}{dt}(f\circ\phi_t^H)=
		\{f,H\}\circ\phi_t^H$; la dinamica classica è codificata nel Lie-bracket con $H$.
	\end{itemize}
	
	% ----------------------------------------------------------
	\subsection{Azione Simplettica e Teorema di Noether}
	% ----------------------------------------------------------
	
	\fonti{Abraham \& Marsden, \textit{Foundations of Mechanics};
		Moretti, \textit{Spectral Theory and Quantum Mechanics}.}
	
	\begin{itemize}
		\item \textbf{Gruppo di Lie $G$:} algebra di Lie $\mathfrak{g}=T_eG$; mappa
		esponenziale; esempi $U(1)$, $SU(2)$, $U(n)$.
		\item \textbf{Azione simplettica:} $\Phi:G\times M\to M$ con $\Phi_g^*\omega=\omega$;
		campo generatore infinitesimale $\xi_M$.
		\item \textbf{Mappa momento:} $\mu:M\to\mathfrak{g}^*$ definita da
		$d\langle\mu,\xi\rangle=\iota_{\xi_M}\omega$. \emph{(Per $M$ non compatta,
			l'esistenza globale richiede la completezza del flusso di $\xi_M$; automatica
			se $M$ è compatta.)}
		\item \textbf{Teorema di Noether geometrico:} se $H$ è $G$-invariante,
		\[
		\frac{d}{dt}\langle\mu\circ\phi_t^H,\xi\rangle = \{\langle\mu,\xi\rangle,H\} = 0.
		\]
		Simmetria $\Leftrightarrow$ conservazione, mediata dal Lie-bracket. Questa
		dimostrazione usa solo la struttura simplettica e l'identità di Jacobi --- si
		applicherà identicamente a $\mathbb{P}\mathcal{H}$ nel Cap.~3.
	\end{itemize}
	
	% ----------------------------------------------------------
	\subsection{$C_0(M)$ come $C^*$-algebra e Teorema di Gelfand}
	% ----------------------------------------------------------
	
	\fonti{Bratteli \& Robinson, \textit{Operator Algebras and Quantum Statistical
			Mechanics}, Vol.~I; Strocchi, \textit{An Introduction to the Mathematical
			Structure of Quantum Mechanics}.}
	
	\nota{Perché $C_0(M)$ e non $C^\infty(M)$? La struttura di $C^*$-algebra richiede
		completezza in norma. $C^\infty(M)$ non ammette una norma di Banach naturale.
		$C_0(M)$ --- funzioni continue che si annullano all'infinito, con norma sup ---
		è lo spazio di completamento naturale su cui la condizione $\|f^*f\|=\|f\|^2$
		è soddisfatta per costruzione. La subalgebra densa su cui la struttura di Poisson
		è effettivamente definita è $C^\infty_c(M)\subset C_0(M)$ --- funzioni lisce a
		supporto compatto, che si annullano all'infinito e sono differenziabili.
		Per $M$ compatta, $C_0(M)=C(M)$ e $C^\infty_c(M)=C^\infty(M)$.}
	
	\begin{itemize}
		\item \textbf{$C_0(M)$ come $C^*$-algebra commutativa:} norma sup, involuzione per
		coniugazione complessa, condizione C* soddisfatta per costruzione.
		\item \textbf{Osservabili classiche come elementi autoaggiunti:} $f=\bar f$;
		spettro reale.
		\item \textbf{Stato come funzionale positivo normalizzato:} $\omega:C_0(M)\to
		\mathbb{C}$ con $\omega(\bar ff)\geq 0$ e $\omega(1)=1$; valore di aspettazione
		come dato fisico primitivo. Approccio parallelo al \S2.1.
		\item \textbf{Teorema di Gelfand:} ogni $C^*$-algebra commutativa è isometricamente
		isomorfa a $C_0(X)$ per qualche spazio localmente compatto di Hausdorff $X$
		\emph{(per $M$ compatta: isomorfa a $C(X)$)}. La MR non \emph{usa}
		$C^*$-algebre --- la MR \emph{è} la teoria delle $C^*$-algebre commutative.
		\item \textbf{Conclusione:} il Lie-bracket classico è struttura interna di
		$C^\infty_c(M)\subset C_0(M)$, non un assioma aggiunto. Il Cap.~2 mostrerà che
		lo stesso vale per il commutatore in $B(\mathcal{H})$.
	\end{itemize}
	
	% ----------------------------------------------------------
	\subsection{Qubit I --- $S^2$ come Spazio delle Fasi dello Spin Classico}
	% ----------------------------------------------------------
	
	\fonti{Abraham \& Marsden, \textit{Foundations of Mechanics};
		Perelomov, \textit{Generalized Coherent States}.}
	
	\begin{itemize}
		\item \textbf{Spazio delle fasi ridotto $S^2$} \emph{(compatto):} forma simplettica
		$\omega=\sin\theta\,d\theta\wedge d\phi$. Su $S^2$ il flusso hamiltoniano è
		automaticamente completo e la mappa momento è globalmente definita.
		\item \textbf{Parentesi di Poisson su $S^2$:} $\{x_i,x_j\}=\varepsilon_{ijk}x_k$;
		algebra $\mathfrak{su}(2)$ realizzata come parentesi di Poisson.
		\item \textbf{$C(S^2)$ come $C^*$-algebra commutativa} (caso compatto: $C_0=C$);
		stato come misura di probabilità su $S^2$.
		\item \textbf{Anticipazione:} nel Cap.~2 il commutatore $[\sigma_i,\sigma_j]=
		2i\varepsilon_{ijk}\sigma_k$ su $M_2(\mathbb{C})$ è la stessa algebra di Lie
		$\mathfrak{su}(2)$ --- non per analogia, ma perché entrambe sono rappresentazioni
		della struttura di \S1.1.
	\end{itemize}
	
	% ============================================================
	\newpage
	\part*{Parte II --- Il Lato Quantistico}
	\addcontentsline{toc}{part}{Parte II: Il Lato Quantistico}
	% ============================================================
	
	\nota{La Parte~II costruisce la MQ senza postulare $\mathcal{H}$. Il percorso è:
		osservabili fisiche $\to$ $C^*$-algebra astratta $\to$ GNS $\to$ $\mathcal{H}$
		come conseguenza. L'ostruzione di Wintner motiva l'algebra di Weyl. Il capitolo
		si chiude aprendo la domanda a cui il Cap.~3 risponde.}
	
	% ============================================================
	\section{Costruzione Algebrica della Meccanica Quantistica}
	% ============================================================
	
	\nota{Struttura parallela al Cap.~1: come $C_0(M)$ è riconosciuta come $C^*$-algebra
		commutativa da Gelfand, $B(\mathcal{H})$ è riconosciuta come $C^*$-algebra non
		commutativa da Gelfand--Naimark. Il Lie-bracket $[a,b]=ab-ba$ è struttura interna,
		non un assioma. La definizione di $C^*$-algebra e di stato non vengono ripetute ---
		sono già in \S1.5.}
	
	% ----------------------------------------------------------
	\subsection{Approccio Operazionale di Strocchi}
	% ----------------------------------------------------------
	
	\fonti{Strocchi, \textit{An Introduction to the Mathematical Structure of Quantum
			Mechanics}.}
	
	\begin{itemize}
		\item \textbf{Impostazione:} preparazioni e misure come punto di partenza; nessuno
		spazio di Hilbert assunto. La struttura $C^*$-algebra è già nota dal \S1.5; ora
		l'algebra è non commutativa.
		\item \textbf{Il Lie-bracket $[a,b]=ab-ba$ come struttura interna:} antisimmetria
		e identità di Jacobi per costruzione algebrica pura. \emph{Parallelo diretto con
			\S1.3: anche lì il Lie-bracket era struttura interna, non un assioma.}
		\item \textbf{Stato come funzionale positivo normalizzato:} approccio parallelo
		a \S1.5.
	\end{itemize}
	
	% ----------------------------------------------------------
	\subsection{Teorema GNS, Gelfand--Naimark e Complementi Tecnici}
	% ----------------------------------------------------------
	
	\fonti{Reed \& Simon, \textit{Methods of Modern Mathematical Physics}, Vol.~I;
		Moretti, \textit{Spectral Theory and Quantum Mechanics}; Strocchi;
		Bratteli \& Robinson, Vol.~I.}
	
	\begin{itemize}
		\item \textbf{Prodotto interno GNS:} dato $\omega$, $\langle a,b\rangle_\omega=
		\omega(a^*b)$; ideale di Gelfand $\mathcal{I}_\omega=\{a:\omega(a^*a)=0\}$.
		\item \textbf{Spazio GNS $\mathcal{H}_\omega$:} completamento di
		$\mathcal{A}/\mathcal{I}_\omega$; rappresentazione $\pi_\omega:\mathcal{A}\to
		B(\mathcal{H}_\omega)$; vettore ciclico $\Omega_\omega=[I]$.
		\item \textbf{Teorema GNS:} ogni stato determina un'unica rappresentazione
		$(\mathcal{H}_\omega,\pi_\omega,\Omega_\omega)$ a meno di equivalenza unitaria;
		$\mathcal{H}$ \emph{emerge} dalla struttura algebrica. Dimostrazione dell'unicità
		via costruzione esplicita dell'isomorfismo unitario.
		\item \textbf{Teorema di Gelfand--Naimark:} ogni $C^*$-algebra astratta si
		rappresenta isometricamente su $B(\mathcal{H})$; lo spazio di Hilbert è il
		supporto naturale di qualsiasi $C^*$-algebra.
		\item \textbf{Complemento tecnico:} costruzione dell'algebra di Weyl per $n>1$
		(necessaria per il \S2.3 e il Cap.~4); si veda anche l'Appendice~A.
	\end{itemize}
	
	% ----------------------------------------------------------
	\subsection{L'Ostruzione di Wintner e l'Algebra di Weyl}
	% ----------------------------------------------------------
	
	\fonti{Wintner, \textit{The unboundedness of quantum-mechanical matrices}
		(Phys.\ Rev.\ 71, 1947); Strocchi;
		Bratteli \& Robinson, Vol.~II; Moretti.}
	
	\nota{Prima di Wintner entrano i prerequisiti di teoria spettrale strettamente
		necessari, integrati come paragrafo organico alla motivazione.}
	
	\begin{itemize}
		\item \textbf{Prerequisito di teoria spettrale:} operatori simmetrici vs
		autoaggiunti su $\mathcal{H}$ --- la distinzione è fondamentale perché
		$\hat q$, $\hat p$ sono non limitati e simmetrici, ma il teorema di Stone
		richiede autoaggiunzione. Teorema spettrale per operatori non limitati (enunciato);
		calcolo funzionale misurabile. Questi strumenti ricompaiono in \S2.3 (Stone)
		e nel Cap.~3.
		\item \textbf{Il problema delle CCR:} si cerca una $C^*$-algebra contenente
		$q,p$ con $[q,p]=i\hbar I$.
		\item \textbf{Teorema di Wintner (1947) --- algebre di Banach:} in qualsiasi
		algebra di Banach unitale non esistono elementi $A,B$ con $AB-BA=I$.
		\begin{itemize}
			\item \textit{Dimostrazione:} per induzione, $AB^n-B^nA=nB^{n-1}$. Quindi
			$n\|B\|^{n-1}\leq\|AB^n-B^nA\|\leq 2\|A\|\|B\|^n$, da cui $n\leq 2\|A\|\|B\|$
			per ogni $n$ --- contraddizione. Questa dimostrazione vale per qualsiasi
			algebra di Banach, senza uso della traccia e senza distinzione di dimensione.
		\end{itemize}
		\item \textbf{Conseguenza:} le CCR $[q,p]=i\hbar I$ non possono essere soddisfatte
		da operatori limitati in nessuna algebra di Banach unitale. Bisogna riformulare.
		\item \textbf{Operatori di Weyl:} $W(\alpha)=e^{i(s\hat q+t\hat p)}$,
		$\alpha=(s,t)\in\mathbb{R}^2$; unitari (limitati) anche se $\hat q$, $\hat p$
		non lo sono.
		\item \textbf{Relazioni di Weyl:} $W(\alpha)W(\beta)=e^{\frac{i}{2}\omega(\alpha,\beta)}
		W(\alpha+\beta)$; CCR in forma esponenziale, senza problemi di dominio.
		\item \textbf{$C^*$-algebra di Weyl $\mathrm{CCR}(\mathbb{R}^2)$:} completamento
		in norma $C^*$; unicità della norma ($\|W(\alpha)\|=1$ dalla condizione C*).
		\item \textbf{Teorema di Stone} \emph{(marcato: riusato nel Cap.~3):} ogni gruppo
		unitario fortemente continuo $\{U(t)\}_{t\in\mathbb{R}}$ su $\mathcal{H}$ ammette
		un unico generatore autoaggiunto $H$ con $U(t)=e^{-iHt/\hbar}$; conversamente,
		ogni operatore autoaggiunto genera tale gruppo.
		\item \textbf{Teorema di Stone--von Neumann:} ogni rappresentazione irriducibile
		fortemente continua di $\mathrm{CCR}(\mathbb{R}^{2n})$ è unitariamente equivalente
		alla rappresentazione di Schrödinger su $L^2(\mathbb{R}^n)$; unicità cinematica
		per sistemi a finiti gradi di libertà \emph{(fallisce in dimensione infinita)}.
	\end{itemize}
	
	% ----------------------------------------------------------
	\subsection{Il Ponte verso la Geometria degli Stati}
	% ----------------------------------------------------------
	
	\fonti{Ashtekar \& Schilling, \textit{Geometrical Formulation of Quantum Mechanics}
		(gr-qc/9706069); Landsman, \textit{Mathematical Topics Between Classical and
			Quantum Mechanics}.}
	
	\nota{Questo paragrafo è la cerniera tra il Cap.~2 e il Cap.~3. Non introduce nuovi
		strumenti tecnici: formula la domanda a cui il Cap.~3 risponde.}
	
	\begin{itemize}
		\item \textbf{La simmetria algebrica:} MR e MQ sono entrambe $C^*$-algebre con
		il Lie-bracket come struttura interna. Nel caso classico quel bracket è il motore
		geometrico che genera il flusso hamiltoniano su $(M,\omega)$, definisce la mappa
		momento, rende Noether un teorema sulla varietà.
		\item \textbf{La domanda:} dato che il Lie-bracket della MQ è la stessa struttura
		astratta, esiste uno spazio in cui le strutture geometriche della MR --- forma
		simplettica, campi hamiltoniani, mappa momento --- siano altrettanto letteralmente
		presenti nel caso quantistico?
		\item \textbf{La traccia della risposta:} il GNS ha prodotto $\mathcal{H}$; Stone
		dice che ogni osservabile autoaggiunto genera un flusso unitario su $\mathcal{H}$.
		Il quoziente $\mathbb{P}\mathcal{H}$ rimuove la ridondanza di fase. La domanda
		diventa: \emph{$\mathbb{P}\mathcal{H}$ è una varietà simplettica? Il flusso
			unitario discende a flusso hamiltoniano su $\mathbb{P}\mathcal{H}$?} Il Cap.~3
		risponde affermativamente.
	\end{itemize}
	
	% ----------------------------------------------------------
	\subsection{Qubit II --- $M_2(\mathbb{C})$ come $C^*$-Algebra}
	% ----------------------------------------------------------
	
	\fonti{Bratteli \& Robinson, Vol.~I; Landsman.}
	
	\begin{itemize}
		\item \textbf{Algebra degli osservabili:} $\mathcal{A}=M_2(\mathbb{C})$; matrici
		hermitiane come osservabili; $[\sigma_i,\sigma_j]=2i\varepsilon_{ijk}\sigma_k$.
		\item \textbf{Stati come matrici densità:} $\rho\geq 0$, $\mathrm{tr}(\rho)=1$;
		stati puri ($\rho^2=\rho$) e misti ($\mathrm{tr}(\rho^2)<1$).
		\item \textbf{GNS esplicito} per lo stato traciale $\omega(a)=\frac{1}{2}
		\mathrm{tr}(a)$: spazio $\mathbb{C}^2$, vettore ciclico $\Omega=\frac{1}{\sqrt{2}}
		(1,1)^T$. Calcolo esplicito.
		\item \textbf{Confronto con Qubit I:} $[\sigma_i,\sigma_j]=2i\varepsilon_{ijk}
		\sigma_k$ in $M_2(\mathbb{C})$ e $\{x_i,x_j\}=\varepsilon_{ijk}x_k$ su $S^2$
		--- stessa algebra di Lie $\mathfrak{su}(2)$, due rappresentazioni distinte.
	\end{itemize}
	
	% ============================================================
	\newpage
	\part*{Parte III --- Le Strutture Classiche nello Spazio degli Stati Quantistico}
	\addcontentsline{toc}{part}{Parte III: Le Strutture Classiche nello Spazio degli
		Stati Quantistico}
	% ============================================================
	
	\nota{La Parte~III offre la risposta top-down alla domanda del \S2.4. Lo spazio degli
		stati $\mathbb{P}\mathcal{H}$ è una varietà simplettica kähleriana. La struttura
		non è assunta: emerge dalla decomposizione del prodotto scalare. Il teorema di
		Stone, enunciato in \S2.3, viene qui riusato in modo esplicito e preciso.}
	
	% ============================================================
	\section{Geometria di $\mathbb{P}\mathcal{H}$ e Recupero delle Strutture Classiche}
	% ============================================================
	
	\nota{Scopo: mostrare che le strutture del Cap.~1 --- forma simplettica, campi
		hamiltoniani, mappa momento, Noether --- sono letteralmente presenti in
		$\mathbb{P}\mathcal{H}$ come conseguenze della geometria intrinseca.}
	
	% ----------------------------------------------------------
	\subsection{Lo Spazio Proiettivo $\mathbb{P}\mathcal{H}$ come Varietà}
	% ----------------------------------------------------------
	
	\fonti{Ashtekar \& Schilling, \textit{Geometrical Formulation of Quantum Mechanics};
		Cirelli, Lanzavecchia \& Mania, \textit{Quantum Mechanics as an
			Infinite-Dimensional Hamiltonian System} (J.\ Math.\ Phys.\ 31, 1990).}
	
	\begin{itemize}
		\item \textbf{Ridondanza di fase:} $|\psi\rangle$ e $e^{i\theta}|\psi\rangle$
		danno gli stessi valori di aspettazione; $\mathbb{P}\mathcal{H}=
		(\mathcal{H}\setminus\{0\})/{\sim}$.
		\item \textbf{Struttura differenziabile reale} di dimensione $2(\dim\mathcal{H}-1)$.
		Per $\mathcal{H}$ finito-dimensionale ($\dim\mathcal{H}=n$):
		$\mathbb{P}\mathcal{H}\cong\mathbb{CP}^{n-1}$, \emph{compatta}. Per $\mathcal{H}$
		infinito-dimensionale separabile: varietà di Hilbert non compatta; le proprietà
		locali (forma simplettica, campi hamiltoniani) restano valide, ma l'esistenza
		globale del flusso richiede che $\hat H$ sia essenzialmente autoaggiunto ---
		condizione garantita dal teorema spettrale di \S2.3.
	\end{itemize}
	
	% ----------------------------------------------------------
	\subsection{Decomposizione del Prodotto Scalare e Struttura di Kähler}
	% ----------------------------------------------------------
	
	\fonti{Ashtekar \& Schilling; Cirelli, Lanzavecchia \& Mania.}
	
	\nota{La definizione di varietà kähleriana come terna $(g,J,\omega)$ con $\omega$
		chiusa entra qui in una riga, al momento in cui serve. La struttura complessa
		integrabile (Nijenhuis, Newlander--Nirenberg) non è necessaria: l'integrabilità
		si verifica direttamente via potenziale di Kähler.}
	
	\begin{itemize}
		\item \textbf{Decomposizione fondamentale:} su $\mathcal{H}$ come spazio
		vettoriale reale,
		\[
		\langle\psi|\phi\rangle = g(\psi,\phi) + i\,\omega(\psi,\phi)
		\]
		dove $g=\mathrm{Re}\langle\cdot|\cdot\rangle$ è una metrica riemanniana e
		$\omega=\mathrm{Im}\langle\cdot|\cdot\rangle$ è una 2-forma antisimmetrica.
		\item \textbf{Non degenerazione di $\omega$:} $\omega$ è simplettica.
		\item \textbf{Struttura complessa $J$:} $\psi\mapsto i\psi$; $J^2=-\mathrm{Id}$;
		compatibilità $\omega(\psi,\phi)=g(J\psi,\phi)$.
		\item \textbf{Discesa a $\mathbb{P}\mathcal{H}$:} la terna $(g,J,\omega)$ discende
		al quoziente e soddisfa le condizioni di compatibilità kähleriana.
		\item \textbf{Potenziale di Kähler e chiusura di $\omega$:} in carta locale
		$U_0$ con coordinate $z_j=\psi_j/\psi_0$, il potenziale $K=\log(1+\sum_j|z_j|^2)$
		genera $\omega=i\partial\bar\partial K$; chiusura $d\omega=0$ dalla nilpotenza
		$\partial^2=0$.
		\item \textbf{Metrica di Fubini--Study:} la componente riemanniana $g$ è la
		metrica di Fubini--Study.
	\end{itemize}
	
	% ----------------------------------------------------------
	\subsection{Equazione di Schrödinger come Flusso Hamiltoniano}
	% ----------------------------------------------------------
	
	\fonti{Ashtekar \& Schilling, \textit{Geometrical Formulation of Quantum Mechanics},
		Prop.~3.1.}
	
	\begin{itemize}
		\item \textbf{Osservabile come funzione liscia:} dato $\hat A$ autoaggiunto,
		$f_A([\psi])=\langle\psi|\hat A|\psi\rangle/\langle\psi|\psi\rangle$ è liscia
		su $\mathbb{P}\mathcal{H}$; analogo diretto dell'osservabile classica di \S1.3.
		\item \textbf{Campo hamiltoniano:} $X_{f_H}$ da $\iota_{X_{f_H}}\omega=df_H$.
		\item \textbf{Stone su $\mathcal{H}$, flusso su $\mathbb{P}\mathcal{H}$:}
		il teorema di Stone (\S2.3) produce il flusso unitario $U(t)=e^{-i\hat Ht/\hbar}$
		su $\mathcal{H}$. Poiché $U(t)$ è lineare e commuta con la moltiplicazione per
		scalari, esso \emph{discende} al quoziente: $\Phi_t:[\psi]\mapsto[e^{-i\hat Ht/\hbar}
		\psi]$. Il flusso disceso coincide con il flusso del campo hamiltoniano $X_{f_H}$
		rispetto a $\omega$ (Ashtekar--Schilling, Prop.~3.1). \emph{Stone lavora su
			$\mathcal{H}$; il flusso scende su $\mathbb{P}\mathcal{H}$ e lì coincide con
			il flusso simplettico. Stone non viene applicato direttamente a
			$\mathbb{P}\mathcal{H}$.}
		\item \textbf{Equazione di Schrödinger:} le orbite di $\Phi_t$ soddisfano
		$i\hbar\partial_t|\psi\rangle=\hat H|\psi\rangle$ --- non un assioma indipendente,
		ma la conseguenza del flusso unitario ridisceso.
		\item \textbf{Identità geometrica fondamentale:}
		\[
		\{f_A,f_B\}_\omega = \omega(X_{f_A},X_{f_B}) = \frac{\langle\psi|[\hat A,\hat B]|\psi\rangle}{i\hbar \langle\psi | \psi \rangle};
		\]
		le parentesi di Poisson su $\mathbb{P}\mathcal{H}$ \emph{sono} i commutatori
		quantistici.
		\item \textbf{Parallelo con il Cap.~1:} stessa struttura, stesso teorema: ogni
		operatore autoaggiunto genera un flusso unitario su $\mathcal{H}$ che discende
		a flusso hamiltoniano su $\mathbb{P}\mathcal{H}$.
	\end{itemize}
	
	% ----------------------------------------------------------
	\subsection{Teorema di Noether Quantistico}
	% ----------------------------------------------------------
	
	\fonti{Cirelli, Lanzavecchia \& Mania (J.\ Math.\ Phys.\ 31, 1990);
		Ashtekar \& Schilling.}
	
	\begin{itemize}
		\item \textbf{Enunciato:} se $[\hat H,\hat A]=0$, allora $\frac{d}{dt}
		\langle\hat A\rangle_\psi=0$ lungo ogni soluzione dell'equazione di Schrödinger.
		\item \textbf{Dimostrazione geometrica:} $[\hat H,\hat A]=0$ implica
		$\{f_H,f_A\}_\omega=0$ per l'identità di \S3.3; per il teorema di Noether di
		\S1.4 applicato a $(\mathbb{P}\mathcal{H},\omega)$, $f_A$ è conservata lungo
		il flusso di $X_{f_H}$. La conservazione segue per identità di Jacobi, garantita
		da $d\omega=0$ su $\mathbb{P}\mathcal{H}$ (\S3.2). \emph{Classico e quantistico
			usano la stessa dimostrazione --- perché usano la stessa struttura.}
		\item \textbf{Mappa momento quantistica:} $\mu_Q:\mathbb{P}\mathcal{H}\to
		\mathfrak{g}^*$ definita da $\langle\mu_Q([\psi]),\xi\rangle=\langle\psi|
		\hat T_\xi|\psi\rangle$; controparte esatta della mappa momento classica di \S1.4.
	\end{itemize}
	
	% ----------------------------------------------------------
	\subsection{Relazioni di Indeterminazione come Corollario}
	% ----------------------------------------------------------
	
	\fonti{Reed \& Simon, Vol.~I; Moretti, \textit{Spectral Theory and Quantum Mechanics}.}
	
	\begin{itemize}
		\item \textbf{Origine dalla struttura simplettica:} Cauchy--Schwarz su $\mathcal{H}$
		dà
		\[
		\Delta_\psi A\cdot\Delta_\psi B \;\geq\;
		\tfrac{1}{2}|\langle\psi|[\hat A,\hat B]|\psi\rangle|;
		\]
		tramite l'identità di \S3.3 il membro destro è $\tfrac{\hbar}{2}|\{f_A,f_B\}_\omega|$.
		\item \textbf{Natura del risultato:} struttura intrinseca della geometria kähleriana
		di $\mathbb{P}\mathcal{H}$, priva di analogo classico diretto. Alcune strutture
		della MR riemergono identiche (Noether, flusso); le relazioni di indeterminazione
		sono genuinamente quantistiche.
	\end{itemize}
	
	% ----------------------------------------------------------
	\subsection{Qubit III --- $\mathbb{P}\mathbb{C}^1 \cong S^2$: Chiusura Geometrica}
	% ----------------------------------------------------------
	
	\fonti{Ashtekar \& Schilling, \textit{Geometrical Formulation of Quantum Mechanics};
		Cirelli, Lanzavecchia \& Mania.}
	
	\begin{itemize}
		\item \textbf{Lo spazio degli stati puri del qubit:} per $\mathcal{H}=\mathbb{C}^2$,
		la costruzione del \S3.1 dà $\mathbb{P}\mathcal{H}=\mathbb{P}\mathbb{C}^1$.
		Topologicamente $\mathbb{P}\mathbb{C}^1\cong S^2$; esplicitamente, la carta
		$U_0=\{[\psi_0,\psi_1]:\psi_0\neq 0\}$ con coordinata $z=\psi_1/\psi_0\in\mathbb{C}$
		copre $S^2$ privata del polo nord tramite la proiezione stereografica.
		
		\item \textbf{La forma simplettica è quella del Qubit I:}
		il potenziale di Kähler del \S3.2 su $\mathbb{P}\mathbb{C}^1$ è
		$K=\log(1+|z|^2)$, e la forma risultante
		$\omega_{FS}=i\partial\bar\partial K=\frac{i\,dz\wedge d\bar z}{(1+|z|^2)^2}$
		in coordinate sferiche $z=\tan(\theta/2)e^{i\phi}$ diventa
		\[
		\omega_{FS} = \tfrac{1}{2}\sin\theta\;d\theta\wedge d\phi.
		\]
		Questa è esattamente la forma simplettica di $S^2$ introdotta nel \S1.6 (Qubit I),
		a meno del fattore di normalizzazione per il raggio dello spin. Lo spazio degli
		stati puri del qubit e lo spazio delle fasi dello spin classico sono la
		\emph{stessa varietà simplettica} --- non un'analogia, un'identità.
		
		\item \textbf{La precessione di Larmor come flusso hamiltoniano su $S^2$:}
		l'hamiltoniana $\hat H=-\frac{\gamma B}{2}\sigma_z$ è autoaggiunta su
		$\mathbb{C}^2$; la funzione osservabile associata è
		$f_H([\psi])=\langle\psi|\hat H|\psi\rangle/\|\psi\|^2$.
		Per il \S3.3, il flusso unitario $U(t)=e^{-i\hat Ht/\hbar}$ discende a un
		flusso hamiltoniano $\Phi_t$ su $(\mathbb{P}\mathbb{C}^1,\omega_{FS})\cong(S^2,\omega)$.
		In coordinate sferiche $\Phi_t$ è la rotazione $\phi\mapsto\phi+\gamma Bt$,
		$\theta\mapsto\theta$ --- precisamente la precessione di Larmor del Qubit I
		come flusso simplettico su $S^2$.
		
		\item \textbf{Chiusura del cerchio geometrico:} il sistema fisico è uno solo.
		Visto come spazio delle fasi classico (Qubit I), vive su $(S^2,\omega)$.
		Visto come spazio degli stati quantistici puri (Qubit II--III), vive su
		$(\mathbb{P}\mathbb{C}^1,\omega_{FS})$. Le due descrizioni usano la stessa
		varietà simplettica con la stessa forma simplettica. La struttura geometrica
		con cui la tesi apre al \S1.2 è la stessa struttura che governa lo spazio degli
		stati del sistema quantistico più semplice.
	\end{itemize}
	
	% ============================================================
	\newpage
	\part*{Parte IV --- Strict Deformation Quantization}
	\addcontentsline{toc}{part}{Parte IV: Strict Deformation Quantization}
	% ============================================================
	
	\nota{La Parte~IV risponde alla domanda complementare: esiste un oggetto matematico
		che connette $C_0(M)$ e $B(\mathcal{H})$ in modo continuo? La risposta di Rieffel
		è affermativa e fornisce una stima normata esplicita tramite la condizione di Dirac.}
	
	% ============================================================
	\section{Quantizzazione per Deformazione (Rieffel)}
	% ============================================================
	
	\nota{Struttura del capitolo: (1) il significato di $\hbar$ come parametro variabile;
		(2) perché la deformazione formale non è sufficiente; (3) la costruzione di Rieffel
		e la condizione di Dirac; (4) la quantizzazione di Weyl come esempio principale;
		(5) l'ostacolo di Groenewold--van Hove come delimitazione della portata.}
	
	% ----------------------------------------------------------
	\subsection{Il Significato di $\hbar$ come Parametro}
	% ----------------------------------------------------------
	
	\fonti{Landsman, \textit{Mathematical Topics Between Classical and Quantum Mechanics}.}
	
	\begin{itemize}
		\item \textbf{(i) $\hbar$ come costante fisica fissata:} regime cinematico
		(Capp.\ 1--2).
		\item \textbf{(ii) $\hbar$ come parametro formale:} variabile nelle serie
		$\mathbb{R}[[\hbar]]$; base della deformazione formale (\S4.2).
		\item \textbf{(iii) $\hbar\in[0,1]$ come indice continuo:} parametrizza le fibre
		di un campo continuo di $C^*$-algebre; il limite classico $\hbar\to 0$ è un
		limite normato (\S\S4.3--4.4).
	\end{itemize}
	
	% ----------------------------------------------------------
	\subsection{Perché la Deformazione Formale Non Basta}
	% ----------------------------------------------------------
	
	\fonti{Landsman; Kontsevich, \textit{Deformation Quantization of Poisson Manifolds}
		(Lett.\ Math.\ Phys.\ 66, 2003).}
	
	\begin{itemize}
		\item \textbf{Star-product formale (BFFLS):} $f\star_\hbar g=\sum_{k=0}^\infty
		\hbar^k B_k(f,g)$; $B_0(f,g)=fg$; $B_1(f,g)-B_1(g,f)=i\{f,g\}$.
		\item \textbf{Problema:} le serie formali non definiscono una $C^*$-algebra;
		nessuna norma di convergenza; il limite $\hbar\to 0$ non ha senso normato.
		\item \textbf{Necessità della strict DQ:} vera $C^*$-algebra per ogni $\hbar$,
		continuità in norma.
	\end{itemize}
	
	% ----------------------------------------------------------
	\subsection{Definizione di Rieffel e Condizione di Dirac}
	% ----------------------------------------------------------
	
	\fonti{Rieffel, \textit{Deformation Quantization for Actions of $\mathbb{R}^d$}
		(Mem.\ AMS, 1993); Landsman.}
	
	\begin{itemize}
		\item \textbf{Definizione (Rieffel, 1993):} una strict deformation quantization
		di $(A_0,\{\cdot,\cdot\})$ è una famiglia $\{A_\hbar\}_{\hbar\in I}$ di
		$C^*$-algebre con $A_0=C_0(M)$ e mappe $Q_\hbar:\widetilde{A}_0\to A_\hbar$
		definite su algebra densa $\widetilde{A}_0\subset A_0$ \emph{(nella pratica:
			$C^\infty_c(M)$ o $\mathcal{S}(\mathbb{R}^{2n})$, dove le parentesi di
			Poisson sono definite)}, tali che:
		\begin{itemize}
			\item[(i)] $Q_\hbar(\bar f)=Q_\hbar(f)^*$;
			\item[(ii)] $\|Q_\hbar(f)\|_{A_\hbar}\to\|f\|_{A_0}$ per $\hbar\to 0$;
			\item[(iii)] \emph{Condizione di Dirac:}
			$\bigl\|\tfrac{1}{i\hbar}[Q_\hbar(f),Q_\hbar(g)]-Q_\hbar(\{f,g\})\bigr\|_{A_\hbar}
			\to 0$ per $\hbar\to 0$.
		\end{itemize}
		\item \textbf{Cosa garantisce Dirac:} non l'uguaglianza tra commutatore e parentesi
		di Poisson, ma la loro differenza piccola in norma $C^*$ per $\hbar$ piccolo.
		Le strutture classiche persistono come \emph{limite normato}.
		\item \textbf{Campo continuo:} classicità e quantisticità connesse da un limite
		asintotico nella topologia della norma --- non da una discontinuità ontologica.
	\end{itemize}
	
	% ----------------------------------------------------------
	\subsection{Esempio Principale: Quantizzazione di Weyl su $\mathbb{R}^2$}
	% ----------------------------------------------------------
	
	\fonti{Landsman, \textit{Mathematical Topics}, Cap.~II.1;
		Hudson, \textit{When is the Wigner quasi-probability density non-negative?}
		(Rep.\ Math.\ Phys.\ 6, 1974).}
	
	\begin{itemize}
		\item \textbf{Setup:} $A_0=C_0(\mathbb{R}^2)$; $A_\hbar=\mathrm{CCR}(\mathbb{R}^2)$;
		$Q_\hbar^W(f)=\frac{1}{2\pi\hbar}\int f(q,p)W(q,p)\,dq\,dp$.
		\item \textbf{Verifica delle tre condizioni di Rieffel:}
		\begin{itemize}
			\item[(i)] compatibilità con l'involuzione: dalle proprietà degli operatori
			di Weyl;
			\item[(ii)] limite della norma: $\|Q_\hbar^W(f)\|\to\|f\|_\infty$ tramite
			stima di Calderón--Vaillancourt;
			\item[(iii)] condizione di Dirac: dal calcolo esplicito del commutatore
			degli operatori di Weyl.
		\end{itemize}
		\item \textbf{Teorema di Hudson:} per uno stato puro $\rho=|\psi\rangle\langle\psi|$,
		la funzione di Wigner $W_\psi\geq 0$ se e solo se $\psi$ è gaussiano.
		\item \textbf{Conclusione:} la corrispondenza di Dirac non è un'analogia ---
		è una disuguaglianza normata.
	\end{itemize}
	
	% ----------------------------------------------------------
	\subsection{Ostacolo di Groenewold--van Hove}
	% ----------------------------------------------------------
	
	\fonti{Moretti, \textit{Spectral Theory and Quantum Mechanics}, p.~605;
		Gotay, \textit{Obstructions to Quantization} (arXiv:math-ph/9809011).}
	
	\begin{itemize}
		\item \textbf{Teorema di Groenewold--van Hove:} non esiste nessuna mappa lineare
		$Q:\mathcal{P}(\mathbb{R}^{2n})\to L(\mathcal{H})$ che soddisfi le regole di
		Dirac esatte sull'algebra dei polinomi; il conflitto emerge sui monomi di grado
		$\leq 4$.
		\item \textbf{Compatibilità con Rieffel:} la condizione di Dirac è \emph{asintotica}
		($\hbar\to 0$), non un'uguaglianza esatta per $\hbar>0$. GvH proibisce l'uguaglianza
		esatta; non proibisce la continuità asintotica.
		\item \textbf{Portata della risposta:} per $\hbar>0$ fissato la quantizzazione
		non può essere un omomorfismo esatto di algebre di Lie su tutto $C^\infty$.
	\end{itemize}
	
	% ----------------------------------------------------------
	\subsection{Qubit IV --- Condizione di Dirac su $\mathfrak{su}(2)$: Verifica Esplicita}
	% ----------------------------------------------------------
	
	\fonti{Landsman, \textit{Mathematical Topics Between Classical and Quantum Mechanics},
		Cap.~II; Rieffel, \textit{Deformation Quantization for Actions of $\mathbb{R}^d$}.}
	
	\begin{itemize}
		\item \textbf{Setup e parametro di deformazione:} si pone $\hbar=1/j$ con
		$j\in\tfrac{1}{2}\mathbb{N}$, $j\to\infty$; la fibra quantistica a spin $j$ è
		$A_j=M_{2j+1}(\mathbb{C})$, con generatori rinormalizzati $\hat x_i=\hat J_i/j$
		(operatori di momento angolare nella rappresentazione di spin $j$). La fibra
		classica è $A_0=C(S^2)$ con coordinate $x_i$ e parentesi
		$\{x_i,x_j\}=\varepsilon_{ijk}x_k$.
		
		\item \textbf{Calcolo del lato quantistico:} le relazioni di commutazione
		$[\hat J_i,\hat J_j]=i\varepsilon_{ijk}\hat J_k$ danno, per i generatori
		rinormalizzati:
		\[
		\frac{1}{i\hbar}[\hat x_i,\hat x_j]
		=\frac{j}{i}[\hat J_i/j,\hat J_j/j]
		=\frac{1}{ij}[\hat J_i,\hat J_j]
		=\frac{\varepsilon_{ijk}\hat J_k}{j}
		=\varepsilon_{ijk}\hat x_k.
		\]
		Dunque $\frac{1}{i\hbar}[\hat x_i,\hat x_j]=\varepsilon_{ijk}\hat x_k$
		\emph{esattamente}, per ogni $j$.
		
		\item \textbf{Verifica della condizione di Dirac:} la mappa di quantizzazione
		$Q_j$ invia $x_i\mapsto\hat x_i$. La condizione di Dirac del \S4.3 richiede
		\[
		\Bigl\|\frac{1}{i\hbar}[Q_j(x_i),Q_j(x_j)] - Q_j(\{x_i,x_j\})\Bigr\|_{A_j}
		\;\to\; 0 \quad\text{per }j\to\infty.
		\]
		Il calcolo mostra che questo scarto è \emph{identicamente zero} per ogni $j$:
		il lato sinistro è $\varepsilon_{ijk}\hat x_k$ e il lato destro è
		$Q_j(\varepsilon_{ijk}x_k)=\varepsilon_{ijk}\hat x_k$. La condizione di Dirac
		è soddisfatta non asintoticamente, ma esattamente --- un caso raro e strutturalmente
		significativo, possibile perché $\mathfrak{su}(2)$ è un'algebra di Lie
		\emph{di dimensione finita} e la quantizzazione è una sua rappresentazione.
		
		\item \textbf{Convergenza degli osservabili nel limite classico:}
		gli stati coerenti di $SU(2)$ $|\theta,\phi;j\rangle$ localizzati nel punto
		$(\theta,\phi)\in S^2$ soddisfano, per $j\to\infty$,
		\[
		\langle\theta,\phi;j|\,\hat x_i\,|\theta,\phi;j\rangle \;\to\; x_i(\theta,\phi),
		\]
		dove $x_i$ sono le coordinate di $S^2$. Il valore di aspettazione dell'operatore
		converge alla funzione classica puntualmente su $S^2$: la fibra quantistica
		$A_j$ converge alla fibra classica $A_0=C(S^2)$ nello spazio degli stati.
		
		\item \textbf{Chiusura del cerchio algebrico:} i quattro Qubit formano un
		percorso chiuso. Il Qubit I pone la struttura di Poisson su $S^2$; il Qubit II
		identifica $M_2(\mathbb{C})$ come $C^*$-algebra non commutativa con lo stesso
		Lie-bracket; il Qubit III mostra che i due spazi degli stati sono la stessa
		varietà simplettica; il Qubit IV mostra che il commutatore normalizzato
		\emph{è} la parentesi di Poisson --- non nel limite, ma per costruzione.
		La domanda guida della tesi trova qui la sua risposta più diretta: le strutture
		non riemergono nel quantistico perché non se ne sono mai andate.
	\end{itemize}
	
	% ============================================================
	\newpage
	\section*{Conclusioni}
	\addcontentsline{toc}{section}{Conclusioni}
	% ============================================================
	
	La domanda corretta non è ``perché le strutture classiche riemergono nel quantistico?''
	ma ``perché ci aspettavamo che non ci fossero?'' Il Cap.~3 mostra che
	$\mathbb{P}\mathcal{H}$ è una varietà simplettica kähleriana e che su di essa vivono,
	nella loro forma originale, tutti gli oggetti del Cap.~1: la forma simplettica emerge
	dalla decomposizione del prodotto scalare, il flusso hamiltoniano discende da Stone,
	Noether segue per identità di Jacobi. Il Cap.~4 lo conferma analiticamente: MR e MQ
	sono fibre dello stesso campo continuo di $C^*$-algebre, e la condizione di Dirac
	garantisce che parentesi di Poisson e commutatore siano la stessa struttura nel senso
	normato per $\hbar\to 0$.
	
	\medskip
	
	\textbf{Delimitazione del dominio.} La risposta è formulata rigorosamente per sistemi
	con un numero finito di gradi di libertà. In QFT il fallimento di Stone--von Neumann
	introduce rappresentazioni unitariamente inequivalenti e fenomeni senza analogo classico:
	il campo continuo di Rieffel non si estende automaticamente a quel contesto. Le
	prospettive aperte riguardano la geometria degli stati misti (metrica di Bures--Uhlmann
	come generalizzazione di Fubini--Study), la decoerenza come meccanismo che sposta il
	sistema verso la fibra $A_0=C_0(M)$, e i sistemi con vincoli tramite riduzione
	simplettica.
	
	% ============================================================
	\newpage
	\section*{Appendici}
	\addcontentsline{toc}{section}{Appendici}
	% ============================================================
	
	\subsection*{Appendice A --- Complementi di $C^*$-Algebre}
	
	\fonti{Bratteli \& Robinson, Vol.~I; Moretti.}
	
	Dimostrazione dell'unicità GNS tramite costruzione esplicita dell'isomorfismo unitario
	tra due rappresentazioni dello stesso stato (citata in \S2.2). Costruzione esplicita
	dell'algebra di Weyl $\mathrm{CCR}(\mathbb{R}^{2n})$ per $n>1$, necessaria per \S2.3
	e \S4.4.
	
	\subsection*{Appendice B --- Tabella di Corrispondenza Strutturale}
	
	\begin{center}
		\renewcommand{\arraystretch}{1.35}
		\begin{tabular}{lll}
			\hline
			\textbf{Struttura} & \textbf{Meccanica Razionale} & \textbf{Meccanica Quantistica} \\
			\hline
			Algebra delle osservabili & $C_0(M)$ (commutativa) & $B(\mathcal{H})$ (non comm.) \\
			Lie-bracket & $\{f,g\}$ & $\tfrac{1}{i\hbar}[A,B]$ \\
			Spazio degli stati & $(M,\omega)$ & $(\mathbb{P}\mathcal{H},\omega)$ \\
			Osservabile come funzione & $f\in C^\infty(M)$ &
			$f_A([\psi])=\langle\psi|\hat A|\psi\rangle/\|\psi\|^2$ \\
			Campo hamiltoniano & $\iota_{X_f}\omega=df$ & $\iota_{X_{f_A}}\omega=df_A$ \\
			Flusso & $\phi_t^H$ su $(M,\omega)$ &
			$[e^{-i\hat Ht/\hbar}\psi]$ su $(\mathbb{P}\mathcal{H},\omega)$ \\
			Conservazione (Noether) & $\{H,f\}=0$ & $[\hat H,\hat A]=0$ \\
			Classificazione & Teorema di Gelfand & Teorema di Gelfand--Naimark \\
			Ponte $\hbar\to 0$ & --- & Condizione di Dirac (Rieffel) \\
			\hline
		\end{tabular}
	\end{center}
	
\end{document}